\documentclass[aspectratio=169]{beamer}
\usepackage{iftex}
\ifPDFTeX
  \usepackage[T1]{fontenc}
  \usepackage{lmodern}
\fi
\usepackage[czech]{babel}
\usepackage{amsmath}
\usepackage{listings}
\usepackage{tikz}
\usetikzlibrary{arrows.meta,positioning,shapes.geometric,babel}

\newcommand{\CVUTTemplatePath}{../../../utils/presentation-template}
\usepackage{\CVUTTemplatePath/beamerthemeCVUT}
\setbeamertemplate{footline}[frame number]
\beamertemplatenavigationsymbolsempty

\newcommand{\LANG}{CZ}
\newcommand{\FRONTPAGE}{dark}
\title{Funkce: definice, volání a rozsah platnosti}
\subtitle{Základy programování\\Matouš Cejnek\\U12110, FS, ČVUT v Praze}
\author{}
\institute{}
\date{}
\renewcommand{\headlineA}{}
\renewcommand{\headlineB}{}
\renewcommand{\headlineC}{}

\definecolor{CodeBlue}{RGB}{38,83,125}
\definecolor{SoftGray}{RGB}{242,244,247}
\definecolor{CodeGreen}{RGB}{38,125,94}
\setbeamerfont{normal text}{size=\large}
\setbeamerfont{frametitle}{size=\LARGE}
\setbeamercolor{block title}{fg=white,bg=CodeBlue}
\setbeamercolor{block body}{fg=black,bg=SoftGray}
\newcommand{\term}[1]{\textcolor{CodeBlue}{\textbf{#1}}}
\lstdefinestyle{python}{language=Python,basicstyle=\ttfamily\large,
  keywordstyle=\color{CodeBlue}\bfseries,stringstyle=\color{CodeGreen},
  commentstyle=\color{gray},showstringspaces=false,columns=fullflexible,keepspaces=true}

\begin{document}

% 0:00--0:01
\begin{frame}
  \clearpage\thispagestyle{empty}\titlepage
\end{frame}

% 0:01--0:03
\begin{frame}{Co si dnes odnesete?}
  \begin{itemize}
    \item Definovat a zavolat vlastní funkci.
    \item Rozlišit parametr, argument a návratovou hodnotu.
    \item Předpovědět pořadí provádění programu s funkcemi.
    \item Vysvětlit lokální a globální rozsah platnosti.
    \item Rozdělit delší program na malé spolupracující části.
  \end{itemize}
\end{frame}

\begin{frame}{Na co navazujeme?}
  \begin{itemize}
    \item Proměnné uchovávají hodnoty různých datových typů.
    \item Podmínky rozhodují a cykly opakují příkazy.
    \item Teď budeme celé kroky algoritmu \term{pojmenovávat a skládat}.
  \end{itemize}
  \begin{block}{Motivace}
    Dobrá funkce převádí jasně určené vstupy na výstup a skrývá podrobnosti výpočtu.
  \end{block}
\end{frame}

% 0:03--0:08
\begin{frame}[fragile]{Když stejný výpočet píšeme opakovaně}
\begin{lstlisting}[style=python]
area_1 = 3.14159 * (d_1 / 2) ** 2
area_2 = 3.14159 * (d_2 / 2) ** 2
area_3 = 3.14159 * (d_3 / 2) ** 2
\end{lstlisting}
  \begin{itemize}
    \item Oprava vzorce se musí provést na více místech.
    \item Záměr kódu není pojmenovaný.
    \item Snadno vzniknou mírně odlišné kopie.
  \end{itemize}
\end{frame}

\begin{frame}{Co je funkce?}
  \begin{block}{Funkce}
    Pojmenovaný blok příkazů, který lze opakovaně zavolat; může přijmout vstupy a vrátit výsledek.
  \end{block}
  \begin{itemize}
    \item \term{Definice} popisuje, co se má při volání stát.
    \item \term{Volání} spustí tělo funkce.
    \item Název vyjadřuje účel: \texttt{circle\_area}, ne \texttt{do\_it}.
  \end{itemize}
\end{frame}

\begin{frame}[fragile]{Definice a volání}
\begin{lstlisting}[style=python]
def circle_area(diameter):
    radius = diameter / 2
    return 3.14159 * radius ** 2

area = circle_area(0.06)
print(area)
\end{lstlisting}
  \begin{itemize}
    \item Za hlavičkou \texttt{def} je dvojtečka.
    \item Odsazení vymezuje tělo funkce.
    \item Závorky ve volání jsou podstatné: \texttt{circle\_area(0.06)}.
  \end{itemize}
\end{frame}

% 0:08--0:15
\begin{frame}[fragile]{Definice sama tělo nespustí}
\begin{lstlisting}[style=python]
def greet(name):
    print("Hello", name)    # runs only when called

print("Before")
greet("Eva")
print("Po")
\end{lstlisting}
  \begin{block}{Výstup}
    \texttt{Před}\quad $\rightarrow$ \quad\texttt{Ahoj Eva}\quad $\rightarrow$ \quad\texttt{Po}
  \end{block}
\end{frame}

\begin{frame}{Jak proběhne volání?}
  \centering
  \begin{tikzpicture}[node distance=6mm,>=Latex,scale=.86,transform shape,
    box/.style={draw,rounded corners,fill=SoftGray,minimum width=34mm,minimum height=9mm,align=center}]
    \node[box] (call) {\texttt{circle\_area(0.06)}};
    \node[box,right=of call] (bind) {\texttt{diameter = 0.06}};
    \node[box,right=of bind] (body) {provedení těla};
    \node[box,right=of body] (ret) {\texttt{return} předá výsledek};
    \draw[->] (call)--(bind); \draw[->] (bind)--(body); \draw[->] (body)--(ret);
  \end{tikzpicture}
  \vspace{4mm}

  Každé volání získá vlastní lokální proměnné. Po návratu program pokračuje za místem volání.
\end{frame}

\begin{frame}{Parametr není argument}
  \begin{description}
    \item[Parametr] jméno vstupu v definici funkce: \texttt{diameter}
    \item[Argument] konkrétní hodnota předaná při volání: \texttt{0.06}
  \end{description}
  \begin{block}{Při volání}
    Argument se přiřadí odpovídajícímu parametru a vznikne lokální proměnná.
  \end{block}
\end{frame}

\begin{frame}[fragile]{Více parametrů: záleží na pořadí}
\begin{lstlisting}[style=python]
def tube_mass(d_out, d_in, length, density):
    area = circle_area(d_out) - circle_area(d_in)
    return area * length * density

mass = tube_mass(0.06, 0.05, 6, 7850)
\end{lstlisting}
  Poziční argumenty se přiřazují zleva doprava.
\end{frame}

\begin{frame}[fragile]{Výchozí a pojmenované argumenty}
\begin{lstlisting}[style=python]
def tube_mass(d_out, d_in, length, density=7850):
    area = circle_area(d_out) - circle_area(d_in)
    return area * length * density

steel = tube_mass(0.06, 0.05, 6)
aluminium = tube_mass(0.06, 0.05, 6,
                      density=2700)
\end{lstlisting}
  \begin{itemize}
    \item Výchozí hodnota dovolí argument vynechat.
    \item Pojmenovaný argument zvyšuje čitelnost volání.
  \end{itemize}
\end{frame}

% 0:15--0:22
\begin{frame}[fragile]{\texttt{print} a \texttt{return} nejsou totéž}
\begin{lstlisting}[style=python]
def double_print(x):
    print(x * 2)

def double_return(x):
    return x * 2

a = double_print(5)   # vytiskne 10, a je None
b = double_return(5)  # nic nevytiskne, b je 10
\end{lstlisting}
\end{frame}

\begin{frame}{Návratová hodnota se dá skládat}
  \begin{itemize}
    \item uložit: \texttt{area = circle\_area(0.06)},
    \item rovnou použít: \texttt{mass = circle\_area(d) * length * density},
    \item předat jiné funkci: \texttt{print(circle\_area(d))},
    \item otestovat bez čtení výpisu.
  \end{itemize}
  \begin{alertblock}{Bez explicitního \texttt{return}}
    Python vrátí speciální hodnotu \texttt{None}.
  \end{alertblock}
\end{frame}

\begin{frame}[fragile]{\texttt{return} funkci okamžitě ukončí}
\begin{lstlisting}[style=python]
def safe_divide(a, b):
    if b == 0:
        return None
    return a / b
    print("Sem se program nikdy nedostane")
\end{lstlisting}
  \begin{itemize}
    \item První provedený \texttt{return} předá výsledek volajícímu.
    \item Kód za ním se v dané větvi neprovede.
  \end{itemize}
\end{frame}

\begin{frame}[fragile]{Více výsledků: ve skutečnosti jedna n-tice}
\begin{lstlisting}[style=python]
def min_max(values):
    smallest = largest = values[0]
    for value in values[1:]:
        if value < smallest:
            smallest = value
        if value > largest:
            largest = value
    return smallest, largest

low, high = min_max([4, 1, 9, 3])
\end{lstlisting}
\end{frame}

% 0:22--0:29
\begin{frame}{Rozsah platnosti (scope)}
  \begin{block}{Lokální jméno}
    Vzniká uvnitř funkce a zvenčí není dostupné.
  \end{block}
  \begin{block}{Globální jméno}
    Je definované na úrovni modulu; funkce ho může číst.
  \end{block}
  Stejné jméno může v různých rozsazích označovat různé proměnné.
\end{frame}

\begin{frame}[fragile]{Lokální proměnná zakryje globální}
\begin{lstlisting}[style=python]
x = 10

def change(x):
    x = x + 1
    print(x)

change(5)  # 6
print(x)   # 10
\end{lstlisting}
  Parametr \texttt{x} je při každém volání nová lokální proměnná.
\end{frame}

\begin{frame}[fragile]{Proč se vyhýbat změnám globálního stavu?}
\begin{lstlisting}[style=python]
total = 0

def add_bad(value):
    global total
    total += value

def add_good(total, value):
    return total + value
\end{lstlisting}
  \begin{itemize}
    \item Skrytá závislost ztěžuje čtení a testování.
    \item Výsledek čistší funkce závisí jen na předaných argumentech.
  \end{itemize}
\end{frame}

\begin{frame}[fragile]{Pozor: měnitelný argument může funkce změnit}
\begin{lstlisting}[style=python]
def add_measurement(values, value):
    values.append(value)

data = [2, 4]
add_measurement(data, 6)
print(data)                 # [2, 4, 6]
\end{lstlisting}
  \begin{alertblock}{Důvod}
    Parametr \texttt{values} odkazuje na stejný seznam jako \texttt{data}. Pokud změnu nechceme, pracujeme s kopií.
  \end{alertblock}
\end{frame}

% 0:29--0:36
\begin{frame}{Jak poznat dobrou funkci?}
  \begin{itemize}
    \item Má jeden jasný účel a výstižné slovesné jméno.
    \item Vstupy dostává parametry; výsledek vrací.
    \item Nespoléhá zbytečně na globální proměnné.
    \item Je krátká natolik, že ji lze přečíst jako jeden krok algoritmu.
    \item Rozhraní je stabilnější než vnitřní postup výpočtu.
  \end{itemize}
\end{frame}

\begin{frame}[fragile]{Funkce skládáme do většího řešení}
\begin{lstlisting}[style=python,basicstyle=\ttfamily\small]
def summary(values):
    low, high = min_max(values)
    return {
        "mean": mean(values),
        "median": median(values),
        "min": low,
        "max": high,
    }
\end{lstlisting}
  \begin{block}{Dekompozice}
    Složitý problém rozdělíme na menší části s jasným rozhraním.
  \end{block}
\end{frame}

\begin{frame}[fragile]{Vstup a výstup oddělíme od výpočtu}
\begin{lstlisting}[style=python,basicstyle=\ttfamily\normalsize]
def celsius_to_fahrenheit(value):
    return value * 9 / 5 + 32

def main():
    value = float(input("Temperature in C: "))
    result = celsius_to_fahrenheit(value)
    print(result)

if __name__ == "__main__":
    main()
\end{lstlisting}
  Výpočet lze znovu použít i ověřit bez klávesnice a obrazovky.
\end{frame}

\begin{frame}[fragile]{Dokumentace funkce: rozhraní před detaily}
\begin{lstlisting}[style=python,basicstyle=\ttfamily\normalsize]
def tube_mass(d_out, d_in, length, density=7850):
    """Return the mass of a hollow tube in kilograms.

    Dimensions are in metres, density in kg/m^3.
    """
    area = circle_area(d_out) - circle_area(d_in)
    return area * length * density
\end{lstlisting}
  \term{Docstring} říká, co funkce očekává a vrací; neopisuje každý řádek.
\end{frame}

% 0:36--0:41
\begin{frame}{Funkce ve vývojovém diagramu}
  \centering
  \vspace{4mm}

  \begin{tikzpicture}[node distance=4mm,>=Latex,scale=.72,transform shape,
    start/.style={draw,ellipse,fill=SoftGray,minimum width=37mm,align=center},
    process/.style={draw,rectangle,fill=SoftGray,minimum width=37mm,minimum height=8mm,align=center},
    decision/.style={draw,diamond,aspect=2,fill=SoftGray,align=center}]
    \node[start] (s) {Funkce \texttt{absolute(x)}};
    \node[decision,below=of s] (d) {$x < 0$?};
    \node[process,below left=7mm and 12mm of d] (neg) {$result=-x$};
    \node[process,below right=7mm and 12mm of d] (pos) {$result=x$};
    \node[start,below=16mm of d] (r) {Návrat \texttt{result}};
    \draw[->] (s)--(d); \draw[->] (d)--node[left]{ano}(neg); \draw[->] (d)--node[right]{ne}(pos);
    \draw[->] (neg) to[out=-90,in=180] (r);
    \draw[->] (pos) to[out=-90,in=0] (r);
  \end{tikzpicture}
  \begin{itemize}
    \item Začátek nese hlavičku funkce, konec její návrat.
    \item Diagram zobrazuje tok uvnitř jednoho volání.
  \end{itemize}
\end{frame}

\begin{frame}[fragile]{Funkce lze předat jako argument}
\begin{lstlisting}[style=python]
def apply_twice(function, value):
    return function(function(value))

def double(x):
    return 2 * x

print(apply_twice(double, 3))  # 12
\end{lstlisting}
  Bez závorek \texttt{double} označuje samotnou funkci; \texttt{double(3)} je její volání.
\end{frame}

\begin{frame}[fragile]{Rekurze: funkce volá sama sebe}
\begin{lstlisting}[style=python]
def factorial(n):
    if n == 0:          # base case
        return 1
    return n * factorial(n - 1)
\end{lstlisting}
  \begin{itemize}
    \item Základní případ zastaví další volání.
    \item Rekurzivní krok převádí problém na menší instanci.
    \item \texttt{factorial(3)} čeká na \texttt{factorial(2)}, ta na \texttt{factorial(1)} atd.
  \end{itemize}
\end{frame}

\begin{frame}{Rekurze není automaticky rychlá}
  \begin{itemize}
    \item Každé volání má vlastní parametry a lokální proměnné.
    \item Naivní Fibonacciho funkce opakovaně počítá stejné podproblémy.
    \item Pomoci může iterace nebo \term{memoizace}: uložení už spočítaných výsledků.
  \end{itemize}
  \begin{alertblock}{Nejdřív správnost}
    Vždy musí existovat základní případ a každý krok se k němu musí přibližovat.
  \end{alertblock}
\end{frame}

% 0:41--0:45
\begin{frame}{Časté chyby}
  \begin{itemize}
    \item Chybí dvojtečka nebo odsazení za \texttt{def}.
    \item Funkce je pouze pojmenovaná, ale nejsou použity závorky pro volání.
    \item Záměna \texttt{print} a \texttt{return}.
    \item Argumenty jsou v jiném pořadí než parametry.
    \item Použití lokální proměnné mimo funkci.
    \item Rekurze nemá dosažitelný základní případ.
  \end{itemize}
\end{frame}

\begin{frame}{Rychlá kontrola porozumění}
  \begin{enumerate}
    \item Co vrátí funkce bez příkazu \texttt{return}?
    \item Jaký je rozdíl mezi parametrem a argumentem?
    \item Proč se globální \texttt{x} nezmění, když má funkce parametr \texttt{x}?
    \item Proč má výpočet hodnotu vracet místo jejího vypsání?
    \item Co musí mít každá rekurzivní funkce?
  \end{enumerate}
\end{frame}

\begin{frame}{Co využijete na cvičení?}
  \begin{itemize}
    \item Skládání funkcí: obsah kruhu $\rightarrow$ průřez $\rightarrow$ hmotnost trubky.
    \item Výchozí a pojmenované argumenty pro hustotu materiálu.
    \item Více návratových hodnot a funkce nad seznamem měření.
    \item Čisté výpočty oddělené od \texttt{input} a \texttt{print} v \texttt{main()}.
    \item Pro pokročilé: funkce jako argument a rekurze.
  \end{itemize}
\end{frame}

\begin{frame}{Shrnutí}
  \begin{itemize}
    \item \texttt{def} vytvoří funkci; závorky ji zavolají.
    \item Parametry tvoří rozhraní, argumenty dodávají konkrétní hodnoty.
    \item \texttt{return} předá hodnotu volajícímu a ukončí funkci.
    \item Lokální proměnné patří jednomu volání; globální stav používejte střídmě.
    \item Malé funkce dovolují program číst, skládat a ověřovat po částech.
  \end{itemize}
\end{frame}

\end{document}
